
          %%%%% ~~~~~~~~~~~~~~~~~~~~ %%%%%

\section{Sampling}

          %%%%% ~~~~~~~~~~~~~~~~~~~~ %%%%%

\subsection{Sampling frame / stratification}
The stratification was by ecoprovince and \textit{stratification} land cover type, i.e., by assigning to 
each sampling unit (grid cell) an ecoprovince as well as a stratification land cover type. 
The ecoprovince assigned to each sampling unit was (of course) the 
ecoprovince that geographically contains that sampling unit,
according to the rasterization of the ecoprovince polygons (vector data)
to the 100 metre $\times$ 100 metre CoE reference geospatial grid \cite{StatCanCoEGridSystem2025}.
On the other hand, each sampling unit was assigned a stratification land cover type by 
resampling by majority rule the LCR map \cite{StatCanCoELCR2025}
(whose underlying grid is the 30 metre $\times$ 30 metre
CoE reference geospatial grid \cite{StatCanCoEGridSystem2025})
to the 100 metre $\times$ 100 metre CoE reference geospatial grid \cite{StatCanCoEGridSystem2025}. 
In what follows, we will refer to this resulting 100 metre $\times$ 100 metre land cover map
as the \textit{stratification map}.

%howpublished= {\url{https://www150.statcan.gc.ca/n1/pub/16-510-x/16-510-x2025002-eng.htm}}
%howpublished= {https://www150.statcan.gc.ca/n1/pub/16-510-x/16-510-x2025002-eng.htm}

\subsection{Sample allocation}
There are 53 ecoprovinces and 10 stratification land cover types (in the simplified LCR map 
typology); hence, there are 530 distinct combinations of ecoprovinces and stratification land 
cover types. However, there were only 423 strata, since certain combinations of 
ecoprovinces and stratification land cover types did not occur in the stratification map. The 
total sample size was 5,011 (number of selected sampling units).
The number of sampling units allocated to each stratum ranged from 3 to 57.
Informed by the results of proof-of-concept and simulation studies
designed to gauge the behaviour of the estimation procedure at the required scale,
sample allocation decisions were made in an expert-driven fashion,
balancing between minimizing the variability of the estimators
for high-priority target variables and ensuring a manageable final sample size. 
 
\subsection{Sampling design}
The sampling design was stratified simple random sampling without replacement. In other 
words, for each stratum (a combination of ecoprovince and stratification land cover type), a 
simple random sample without replacement of fixed size (according to the chosen sample 
allocation) was chosen from the collection of sampling units belonging to that stratum, and 
the sample for each stratum was chosen independently of the sample from every other 
stratum.

\begin{center}
\vskip 0.1cm
\mbox{}
\includegraphics[width=4.25in,clip,trim={20 65 5 75}]{\graphicsDir/official-stratified-initial-sample-100m}
\vskip 0.01cm
\begin{minipage}{4.25in}
{\small An illustration what the spatial distribution of the selected sampling units could look like.
The sampling units are indicated as yellows dots.
The ecoprovince boundaries are indicated in blue.}
\end{minipage}
\end{center}

          %%%%% ~~~~~~~~~~~~~~~~~~~~ %%%%%
